\begin{textsample}{POS Dim 2 – human – Score 6.00 – rj/rj\_ef\_ciencias\_humanas\_his\_2\_p13.txt}  \label{ex:f2_pos_018}
Mais uma vez, enfatizamos que essa relação passado/presente só se faz possível se conseguirmos desenvolver uma “ atitude historiadora ” em cada estudante, se cada indivíduo for \textbf{capaz} de comparar, interpretar, analisar e contextualizar o seu passado, com autonomia de pensamento, e pensamento crítico, para que o conhecimento histórico possa se construir de forma sólida e conectada com a realidade, para que assim, o ensino de história possa estar em consonância com o “ espírito ” da BNCC e da área de \textbf{humanas}, contribuindo, antes de tudo para a formação de seres humanos com valores éticos e dispostos a construir uma \textbf{nova} \textbf{sociedade}.

% matched lemmas: capaz, humano, novo, sociedade
\end{textsample}
